\documentclass[a4paper]{article}
\usepackage{amsmath}
\usepackage{graphicx}
\usepackage{indentfirst}
\usepackage{listings}
\usepackage{matlab-prettifier}
\usepackage{ifthen}

\author{Sergiusz Warga}
\title{Partial Differential Equations with Applications in Physics and Industry}

\begin{document}

\maketitle

\section{The Problem: The Anatomy of a Traffic Jam}

Suppose you are modeling traffic flow on a single-lane, infinitely long highway.
Let $u(x, t)$ represent the normalized \textit{traffic density}, where $u = 0$
is an empty road and $u = 1$ is complete traffic jam.

Observations show that drivers adjust their velocity $v$ based on the density of
cars around them according to the \textbf{Greenshields model}:
\begin{equation}
  v(u) = v_{max}(1 - u),
\end{equation}
where the maximum speed limit is normalized to $v_max = 1$.

At time $t = 0$, the highway experiences a sudden event: cars approaching from
the left ($x < 0$) are flowing at a light density of $x = 0.25$.
However, starting exactly at ($x = 0$), an accident has cause a complete traffic
jam ($u = 1.0$) that stretches indefinitely to the right ($x \geq 0$).

\subsection{Part (a): The Conservation Law}

Assuming no cars enter or leave the highway (\textit{conservations of mass}),
derive the governing non-linear PDE for the traffic density ($u(x, t)$).
Express it in the form $u_t + q(u)_x = 0$, and determine the mathematical
expression for the characteristic wave speed $c(u)$.

\end{document}
